% Vietnamese translation for OI Public Library
% Source-PDF: IOI中国国家候选队论文/2002/方奇.pdf
\documentclass[11pt]{article}
\usepackage{oipl-vn}

\title{Bàn về ứng dụng của điều kiện cần}
\author{Fang Qi}
\date{}

\begin{document}
\maketitle

\origtitle{On the application of necessary conditions}
\sourcepath{IOI China candidate team papers / 2002 / Fang Qi.pdf}

\section*{Tóm tắt}

Điều kiện cần là một quan hệ logic giữa các mệnh đề. Trong quá trình giải
bài, nếu biết dùng tốt điều kiện cần, ta có thể nâng cao hiệu quả thuật toán.
Bài viết trước hết nêu định nghĩa của điều kiện cần, sau đó dựa trên đặc điểm
của nó để phân tích các ứng dụng trong nhiều dạng bài toán, chỉ ra rằng điểm
mấu chốt khi dùng điều kiện cần là ``giảm phần dư thừa, thể hiện bản chất''.
Cuối cùng, bài viết tổng kết một vài kinh nghiệm khi tìm điều kiện cần.

\paragraph{Từ khóa.}
Mệnh đề; điều kiện cần; hiệu quả; độ chính xác.

\section{Định nghĩa điều kiện cần}

Trong toán học, một câu có thể phán định đúng hoặc sai được gọi là một
\term{mệnh đề}. Ta thường dùng các chữ cái Latin thường $p,q,r,s,\ldots$ để
biểu diễn mệnh đề. Nếu từ mệnh đề $p$ có thể suy ra mệnh đề $q$, tức là
``nếu $p$ đúng thì $q$ đúng'', ta viết
\[
  p\Rightarrow q.
\]

Nói chung, nếu đã biết $p\Rightarrow q$, thì $p$ là \term{điều kiện đủ} của
$q$, còn $q$ là \term{điều kiện cần} của $p$. Ví dụ, $x^2>0$ là điều kiện cần
của $x>0$. Ở đây tập nghiệm của $x^2>0$ chứa tập nghiệm của $x>0$:
\[
  \{x\mid x>0\}\subseteq \{x\mid x^2>0\}.
\]
Có thể hình dung bằng một biểu đồ Venn: miền ứng với điều kiện $x>0$ nằm bên
trong miền lớn hơn ứng với điều kiện $x^2>0$.

\section{Ứng dụng của điều kiện cần}

Trong các bài toán tin học, điều kiện cần có phạm vi ứng dụng rất rộng. Bài
viết tập trung vào hai mặt chính: thu hẹp phạm vi cần xét và làm lộ ra bản
chất của bài toán.

\subsection{Thu hẹp phạm vi tìm kiếm}

Phần lớn bài toán yêu cầu tìm các phương án hoặc giá trị cụ thể thỏa một số
điều kiện. Gọi các điều kiện ấy là mệnh đề $p$, và gọi tập nghiệm là $A$.
Khi giải bài, dĩ nhiên ta không thể biết ngay nghiệm chính xác của $A$. Thông
thường ta xuất phát từ điều kiện đã biết $q$, xét trong một phạm vi lớn hơn
$B$, rồi dùng liệt kê, tìm kiếm, quy hoạch động, tham lam hoặc những phương
pháp khác để sàng lọc, cuối cùng tìm ra tập nghiệm $A$ thỏa $p$.

Trong quá trình ấy, $B$ chứa $A$, và mọi nghiệm trong $A$ cũng thỏa điều kiện
$q$. Do đó $p\Rightarrow q$, tức $q$ là điều kiện cần của $p$. Ta bắt đầu từ
điều kiện $q$ tương đối ``lỏng'', rồi tìm trong $B$ các nghiệm thỏa điều kiện
$p$ chặt hơn. Có thể thấy độ lỏng của điều kiện cần ban đầu, tức kích thước
của $B$, ảnh hưởng trực tiếp đến hiệu quả giải bài. Khi $q$ càng gần $p$, tập
$B$ càng gần $A$, tỉ lệ ``vàng'' trong hỗn hợp càng cao, phần ``cát'' cần
loại bỏ càng ít, và hiệu quả tự nhiên càng tốt.

Tuy vậy, điều này phải đặt trên cơ sở khả thi. Nếu điều kiện $q$ khó kiểm tra
hoặc tốn rất nhiều thời gian xử lý, thì dù $q$ có chính xác đến đâu cũng mất
ý nghĩa. Vì vậy, dưới nguyên tắc ``dễ cài đặt'', ta nên cố gắng tìm điều kiện
cần càng chính xác càng tốt để thu hẹp phạm vi giải và tăng tốc độ tìm
nghiệm.

\subsubsection*{Ví dụ 1: cạnh bắt buộc trong ghép hoàn hảo}

\paragraph{Mô tả bài toán.}
Cho đồ thị hai phía $G(V,E)$. Hỏi những cạnh nào là bắt buộc trong mọi ghép
hoàn hảo. Đồ thị có tổng cộng $2n$ đỉnh, với $n\le 100$.

\paragraph{Phân tích.}
Gọi tập các cạnh bắt buộc là $A$. Một cạnh $e$ thuộc $A$ khi và chỉ khi sau
khi xóa $e$, đồ thị con $g=G-e$ không còn ghép hoàn hảo. Đây có thể xem là
điều kiện $p$. Từ đó ta nhận được thuật toán đầu tiên:
\begin{enumerate}
  \item Liệt kê mọi cạnh $e\in E$ của đồ thị $G$.
  \item Xóa tạm $e$ khỏi đồ thị, đặt $g=G-e$.
  \item Tìm ghép cực đại của $g$.
  \item Nếu $g$ không có ghép hoàn hảo thì $e\in A$, ngược lại thì không.
\end{enumerate}

Dùng thuật toán Hungary để tìm ghép cực đại mất $O(n^2)$, còn số cạnh cần
liệt kê có thể là $O(n^2)$. Vì thế độ phức tạp là $O(n^4)$, quá lớn.

Trong thuật toán trên, theo lối nghĩ thông thường, ta vô tình đặt phạm vi ban
đầu $B$ bằng toàn bộ tập cạnh $E$, tức điều kiện cần $q$ chỉ là ``cạnh thuộc
$E$''. Phạm vi tìm kiếm vì vậy bị đặt quá rộng.

Thực ra, nếu $e\in A$, thì $e$ phải thuộc bất kỳ ghép hoàn hảo nào. Nói cách
khác, mọi ghép hoàn hảo đều chứa tất cả các cạnh bắt buộc. Ta có thể cải tiến:
\begin{enumerate}
  \item Trước hết tìm một ghép hoàn hảo $M$ của $G$.
  \item Chỉ liệt kê các cạnh $e\in M$.
  \item Với mỗi cạnh ấy, xóa tạm $e$, tìm ghép cực đại của $G-e$.
  \item Nếu $G-e$ không còn ghép hoàn hảo thì $e\in A$, ngược lại thì không.
\end{enumerate}

Ở đây, phạm vi ban đầu được thu hẹp thành một ghép hoàn hảo, và điều kiện cần
được làm chính xác thành ``cạnh thuộc ghép hoàn hảo $M$''. Độ phức tạp giảm
còn $O(n^3)$.

Ví dụ này cho thấy điều kiện cần rất hữu ích trong các bài toán tồn tại.
Sức mạnh của việc thu hẹp phạm vi còn thể hiện rõ hơn trong tìm kiếm.
Những thuật toán như DFS hay BFS không có heuristic thường có độ phức tạp
hàm mũ; vì vậy, giảm lượng tìm kiếm là điểm then chốt. Dùng điều kiện cần để
cắt bỏ những nhánh không còn hy vọng chính là thao tác thường gọi là
\term{cắt tỉa}.

\subsubsection*{Ví dụ 2: Puzzle, GDOI 1998}

\paragraph{Mô tả bài toán.}
Xem phụ lục về bài \emph{Puzzle}.

\paragraph{Phân tích.}
Đây là một bài tìm kiếm điển hình. Theo đề bài, ta biết:
\begin{enumerate}
  \item cần biến trạng thái ban đầu thành trạng thái đích trong hữu hạn bước,
  hoặc xác định là vô nghiệm;
  \item số bước di chuyển đã biết;
  \item số lần di chuyển theo từng hướng đã biết, nhưng thứ tự chưa biết;
  \item nghiệm phải đồng thời thỏa các ràng buộc ở hai ý trên.
\end{enumerate}

Ta có thể dùng DFS để lần lượt sinh các dãy di chuyển khả dĩ và in ra dãy
thỏa đề. Tìm kiếm toàn bộ như vậy rõ ràng không chấp nhận được về thời gian.
Có thể cắt tỉa từ các mặt sau.

\paragraph{Số lần của từng hướng di chuyển.}
Vì số lần của từng hướng đã biết, nếu trong quá trình tìm kiếm số lần dùng
một hướng nào đó vượt giới hạn thì không cần tiếp tục theo hướng ấy.

\paragraph{Vị trí ô trống và số lần của các hướng.}
Do trạng thái hiện tại, trạng thái đích và số lần còn lại của từng hướng đều
đã biết, ta có thể kiểm tra liệu ô trống hiện tại có thể đi đến vị trí ô trống
đích bằng tập hướng còn lại hay không. Cụ thể, cần thỏa
\[
  x_{\text{đích}}-x_{\text{hiện tại}}
  = \#4_{\text{còn lại}}-\#3_{\text{còn lại}},
\]
và
\[
  y_{\text{đích}}-y_{\text{hiện tại}}
  = \#2_{\text{còn lại}}-\#1_{\text{còn lại}}.
\]
Ở đây các số $1,2,3,4$ lần lượt là bốn thao tác trong đề bài.

\paragraph{Vị trí các ký tự và số lần của các hướng.}
Tương tự, dựa trên trạng thái hiện tại, trạng thái đích và số lần còn lại của
các hướng, ta kiểm tra từng ký tự có còn khả năng đi đến vị trí đích hay
không.

Trong lập luận trên, đặt $p$ là mệnh đề ``dãy thao tác có thể đưa trạng thái
đầu về trạng thái cuối''. Gọi tập các dãy thao tác thỏa $p$ là $A$. Ba phép
cắt tỉa tương ứng với các điều kiện $q_1,q_2,q_3$, và các tập dãy thao tác
thỏa chúng lần lượt là $B_1,B_2,B_3$. Khi đó
\[
  p\Rightarrow q_1,\qquad p\Rightarrow q_2,\qquad p\Rightarrow q_3,
\]
nên
\[
  A\subseteq B_1\cap B_2\cap B_3.
\]
Biểu đồ trong bài gốc minh họa $A$ nằm trong phần giao của ba miền $B_1$,
$B_2$ và $B_3$. Ba điều kiện cần cùng tác dụng làm phạm vi tìm kiếm nhỏ và
chính xác hơn. Cần chú ý rằng ba phép cắt tỉa này đều dễ cài đặt và không tốn
quá nhiều thời gian; nếu không thì sẽ lợi bất cập hại. Sau những tối ưu ấy,
lượng tìm kiếm giảm mạnh và chương trình chạy nhanh hơn nhiều.

\subsubsection*{Ví dụ 3: Musketeers, POI 1999/2001}

\paragraph{Mô tả bài toán.}
$n$ người đứng thành vòng tròn và được đánh số từ $1$ đến $n$. Các trận quyết
đấu diễn ra giữa hai người kề nhau; người thắng ở lại trong vòng, người thua
bị loại. Không biết thứ tự các trận đấu, chỉ biết tổng cộng có $n-1$ trận.
Cho quan hệ thắng thua giữa mọi cặp người, hỏi những ai có thể là người thắng
cuối cùng. $n\le 100$.

\paragraph{Phân tích.}
Bài này tương tự bài \emph{Gộp đá} của NOI 1995. Bắt chước bài đó, ta có một
cách quy hoạch động.

Gọi $E[i,j]$ là kết quả thắng thua giữa $i$ và $j$: nếu $i$ thắng $j$ thì
$E[i,j]=\mathit{true}$, nếu $j$ thắng $i$ thì $E[i,j]=\mathit{false}$.
Gọi $A[i,j,k]$ cho biết người $k$ có thể thắng trong đoạn liên tiếp từ $i$
đến $j$ hay không, trong đó $k$ nằm trong đoạn ấy. Khi có thể thì giá trị là
\textit{true}, ngược lại là \textit{false}. Công thức chuyển trạng thái là
\[
A[i,j,k]=
\begin{cases}
\mathit{true}, & \begin{array}{l}
\text{tồn tại } i\le y\le x<k \text{ sao cho}\\
A[i,x,y]\land A[x+1,j,k]\land E[k,y],\\
\text{hoặc tồn tại } k<u\le v\le j \text{ sao cho}\\
A[u,j,v]\land A[i,u-1,k]\land E[k,v];
\end{array}\\
\mathit{false}, & \text{ngược lại.}
\end{cases}
\]

Thuật toán này có độ phức tạp thời gian $O(n^5)$ và bộ nhớ $O(n^3)$, quá lớn
để dùng trực tiếp.

Chú ý rằng điều kiện cần để người số $x$ thắng trong toàn bộ vòng là người ấy
có thể ``gặp chính mình''. Ta xem vòng như một đoạn thẳng bằng cách tách điểm
$x$ thành hai bản sao ở hai đầu; mọi người ở giữa bị loại, còn $x$ không thua.
Với cách nhìn này, trong một đoạn liên tiếp chỉ cần xét hai người ở hai đầu
có thể thắng để gặp nhau hay không, không cần lưu người thắng ở giữa. Như vậy
trạng thái giảm đi một chiều.

Gọi $\mathit{meet}[i,j]$ cho biết $i$ và $j$ có thể gặp nhau hay không. Công
thức chuyển trạng thái là
\[
\mathit{meet}[i,j]=
\begin{cases}
\mathit{true}, &
\text{tồn tại } i<k<j \text{ sao cho }
\mathit{meet}[i,k]\land \mathit{meet}[k,j]\land (E[i,k]\lor E[j,k]);\\
\mathit{false}, & \text{ngược lại.}
\end{cases}
\]

Độ phức tạp thời gian giảm còn $O(n^3)$, và bộ nhớ giảm còn $O(n^2)$. Trong
bài này, hai cách chọn điều kiện cần về thắng thua dẫn đến hai lối quy hoạch
động khác nhau. Cách sau chính xác hơn, giảm tính toán dư thừa và cải thiện
cả thời gian lẫn bộ nhớ.

Các ví dụ trên phân tích cách dùng điều kiện cần để thu hẹp phạm vi trong
liệt kê, tìm kiếm và quy hoạch động. Ta thấy khi điều kiện cần được chọn
chính xác, khác biệt trước và sau tối ưu là rất rõ rệt.

Nếu ``thu hẹp phạm vi tìm kiếm'' chủ yếu là tối ưu trên thuật toán đã có, thì
phần tiếp theo cho thấy việc chọn đúng điều kiện cần còn có thể làm lộ bản
chất bài toán, từ đó dẫn ta đến thuật toán mới phản ánh tốt hơn các quan hệ
bên trong.

\subsection{Làm lộ bản chất của bài toán}

\subsubsection*{Ví dụ 4: Experiment, GDOI 2001}

\paragraph{Mô tả bài toán.}
Xem phụ lục về bài \emph{Experiment}.

\paragraph{Phân tích.}
Nếu trừu tượng hóa $n_1$ và $n_2$ mảnh thiên thạch thành các đỉnh, còn bỏ qua
các thông tin khác, bài toán trở thành bài tìm ghép tốt nhất trên đồ thị hai
phía. Có thể dùng thuật toán Hungary hoặc luồng chi phí để giải. Độ phức tạp
cỡ $O(n_1^2n_2^2)$; khi $n_1,n_2$ đều đạt giá trị tối đa $500$ thì không thể
chịu được.

Thực ra, nếu chỉ xem đây là bài ghép thì ta mới thấy tính phổ biến của mâu
thuẫn, nhưng chưa thấy tính đặc thù. Suy nghĩ kỹ hơn, trọng số trên các cạnh
của đồ thị hai phía không được gán tùy ý, mà bằng trị tuyệt đối của hiệu khối
lượng giữa hai đỉnh được nối. Đây là khác biệt lớn nhất giữa mô hình của bài
này và đồ thị hai phía tổng quát.

Bài gốc dùng ba hình nhỏ để minh họa các cách nối khi đặt hai dãy khối lượng
trên trục tọa độ. Sau khi sắp xếp khối lượng các mảnh từ sao $A$ và sao $B$,
một điều kiện cần của ghép tối ưu là các cạnh được chọn không cắt nhau. Trường
hợp hai cạnh cắt nhau và không cắt nhau trong một cấu hình đặc biệt có hiệu
quả như nhau, nên ta quy ước chỉ chọn phương án không cắt nhau.

Nói cách khác, giả sử một thí nghiệm chọn mảnh có khối lượng $a_1$ từ sao $A$
và $b_1$ từ sao $B$, còn một thí nghiệm khác chọn $a_2,b_2$. Để tổng hiệu
khối lượng là tốt nhất, không thể tồn tại hai tình huống
\[
  a_1<a_2 \text{ và } b_1>b_2,
  \qquad\text{hoặc}\qquad
  a_1>a_2 \text{ và } b_1<b_2.
\]

Với điều kiện này, bài toán đồng thời thỏa tính không hậu hiệu và nguyên lý
tối ưu, nên có thể giải bằng quy hoạch động. Trước hết sắp xếp hai dãy số từ
nhỏ đến lớn. Giả sử $n_1\le n_2$; nếu không thì đổi vai trò hai dãy. Mảng
$A[i,j]$ biểu diễn hiệu quả tốt nhất khi xét $i$ mảnh đầu của nhóm thứ nhất
và $j$ mảnh đầu của nhóm thứ hai. Công thức chuyển là
\[
A[i,j]=
\begin{cases}
\min\{A[i,j-1],\, A[i-1,j-1]+|w_1[i]-w_2[j]|\}, & j>i,\\
A[i-1,j-1]+|w_1[i]-w_2[j]|, & j=i.
\end{cases}
\]
Trong đó $w_1$ lưu khối lượng các mảnh của nhóm thứ nhất, còn $w_2$ lưu khối
lượng các mảnh của nhóm thứ hai.

Thuật toán mới mất $O(n_1\log n_1+n_2\log n_2)$ để sắp xếp và $O(n_1n_2)$ cho
quy hoạch động, nên độ phức tạp tổng thể là $O(n_1n_2)$. Đây chính là kết quả
của việc điều kiện cần đã chạm đúng điểm cốt lõi của bài toán.

\section{Cách tìm điều kiện cần}

Qua các ví dụ điển hình ở trên, có thể rút ra một vài kinh nghiệm khi tìm
điều kiện cần.

\paragraph{Quy nạp cái chung từ đặc thù.}
Như cách xác định tập $B$ trong ví dụ 1, tính phổ biến thường nằm trong tính
đặc thù và biểu hiện thông qua đặc thù. Điều này cũng phù hợp với thói quen
tư duy của ta: khi gặp một bài khó, ta thường thử bắt đầu từ những trường hợp
đặc biệt đơn giản nhất, dần dần đi sâu hơn, rồi cuối cùng phân tích và quy
nạp ra quy luật tổng quát.

\paragraph{Tìm điều kiện cần từ nhiều mặt.}
Đôi khi, đặc biệt trong thuật toán tìm kiếm, một phép cắt tỉa đơn lẻ là chưa
đủ mạnh. Khi các yếu tố bên trong tác động lẫn nhau và khó nắm bắt toàn cục,
ta có thể thử phân tích bài toán từ nhiều phía để tìm các điều kiện cần, tách
bài toán thành nhiều mặt, rồi tổng hợp các điều kiện cần ấy khi sử dụng. Cách
làm này thường tạo ra hiệu quả cộng hưởng; bài \emph{Puzzle} trong ví dụ 2 là
một trường hợp điển hình.

\paragraph{So sánh với mô hình ban đầu để phân tích giống và khác.}
Bản chất thường hiện ra trong so sánh. Ở bài \emph{Experiment}, chính việc so
sánh với mô hình đồ thị hai phía tổng quát đã làm lộ ra khác biệt căn bản:
điều kiện cần ``các cạnh không cắt nhau''. Từ đó ta được dẫn đến lời giải bằng
quy hoạch động.

Phương pháp là linh hoạt và đa dạng. Khi tìm điều kiện cần, ta nên giữ nguyên
tắc phân tích từng bài toán cụ thể, không bó buộc vào một khuôn mẫu duy nhất.

\section{Kết luận}

Điều kiện cần thực chất là một cơ sở lý thuyết của suy luận logic, đồng thời
cũng là một cách thể hiện quá trình suy nghĩ. Nó có ứng dụng rộng rãi trong
nhiều mặt. Khi giải bài, nếu tìm được điều kiện cần chính xác, ta vừa có thể
làm rõ bản chất của bài toán, vừa có thể thu hẹp phạm vi tìm kiếm, từ đó nâng
cao đáng kể hiệu quả thuật toán.

Bản thân việc tìm điều kiện cần chính xác không có một quy tắc cố định. Điều
đó đòi hỏi ta phải chịu khó suy nghĩ, biết quan sát phát hiện và thường xuyên
tổng kết trong thực hành.

\appendix
\section{Bài Puzzle}

Có một bàn cờ ô vuông $5\times 5$. Trên bàn có $24$ quân cờ khác nhau, lần
lượt được ký hiệu bằng các chữ cái in hoa \texttt{A}, \texttt{B}, \ldots,
\texttt{X}; còn một ô để trống, ký hiệu bằng dấu cách.

Mỗi bước của trò chơi di chuyển quân cờ ở phía trên, phía dưới, bên trái hoặc
bên phải ô trống vào ô trống. Bốn thao tác này lần lượt được ký hiệu là
$1,2,3,4$.

Nếu biết trạng thái ban đầu của bàn cờ và một dãy thao tác hữu hạn theo đúng
thứ tự, ta nhận được duy nhất một trạng thái đích. Chẳng hạn, trong ví dụ của
đề gốc, trạng thái đầu qua dãy thao tác \texttt{144223} sẽ đến trạng thái
đích tương ứng. Tuy nhiên, thứ tự thao tác đúng ban đầu đã bị xáo trộn: nếu
thực hiện dãy bị xáo trộn từ trạng thái đầu thì không nhận được trạng thái
đích. Chỉ có thứ tự bị xáo trộn, còn số lượng từng loại thao tác không đổi.

Nhiệm vụ là tìm lại một dãy thao tác đúng ban đầu. Nếu có nhiều nghiệm, chỉ
cần đưa ra một nghiệm bất kỳ.

\paragraph{Dữ liệu vào.}
Đọc từ tệp \texttt{PUZZLE.DAT} trong thư mục hiện tại. Năm dòng đầu là trạng
thái đầu, mỗi dòng có năm ký tự. Năm dòng tiếp theo là trạng thái đích, mỗi
dòng có năm ký tự. Dòng thứ mười một là số nguyên dương $L$, độ dài dãy thao
tác ($L\le 50$). Dòng thứ mười hai có $L$ ký tự, biểu diễn dãy thao tác đã bị
xáo trộn. Trong mỗi dòng không có ký tự thừa ở đầu hoặc cuối.

\paragraph{Dữ liệu ra.}
Ghi vào tệp \texttt{PUZZLE.OUT}. Nếu có nghiệm, in một dòng gồm $L$ ký tự là
dãy thao tác đúng ban đầu. Nếu vô nghiệm, in \texttt{0}.

\paragraph{Ví dụ.}
\begin{verbatim}
PUZZLE.DAT
TRGSJ
XDOKI
M VLN
WPABE
UQHCF
TRGSJ
XOKLI
MDVBN
WP AE
UQHCF
6
442231

PUZZLE.OUT
144223
\end{verbatim}

\section{Bài Musketeers}

Vào thời vua Louis XIII và vị hồng y Richelieu quyền lực, trong quán trọ Full
Barrel, $n$ chàng lính ngự lâm ăn xong bữa và đang uống rượu. Rượu chưa cạn,
nên họ bắt đầu gây sự; một cuộc ẩu đả say xỉn nổ ra, trong đó mỗi người đều
xúc phạm tất cả những người còn lại.

Một cuộc quyết đấu là không tránh khỏi. Nhưng ai đấu với ai, và theo thứ tự
nào? Họ quyết định đứng thành vòng tròn và rút thăm theo thứ tự. Người được
rút thăm sẽ đấu với người bên phải mình. Kẻ thua rời khỏi trò chơi, chính xác
hơn là thi thể của người ấy được người hầu mang đi. Người kế tiếp đứng cạnh
kẻ thua sẽ trở thành hàng xóm của người thắng.

Nhiều năm sau, khi các sử gia đọc hồi ký của người thắng cuộc, họ nhận ra
rằng kết quả cuối cùng phụ thuộc rất lớn vào thứ tự các trận quyết đấu. Họ
cũng biết từ luyện kiếm rằng ai có thể thắng ai trong từng cặp. Quan hệ
``$A$ thắng $B$'' không có tính bắc cầu: có thể $A$ giỏi hơn $B$, $B$ giỏi hơn
$C$, nhưng $C$ lại giỏi hơn $A$. Vì vậy, trong ba người ấy, trận đầu tiên ảnh
hưởng đến kết quả cuối cùng. Nếu $A$ và $B$ đấu trước thì cuối cùng $C$ thắng;
nếu $B$ và $C$ đấu trước thì cuối cùng $A$ thắng. Các sử gia muốn xác định
những người nào có thể sống sót.

\paragraph{Nhiệm vụ.}
$n$ người đánh số liên tiếp từ $1$ đến $n$ đứng thành vòng tròn. Họ đấu
$n-1$ trận. Ở vòng đầu, một người nào đó, ví dụ số $i$, đấu với người bên phải
mình, tức người số $i+1$ hoặc số $1$ nếu $i=n$. Người thua rời trò chơi và
vòng tròn được khép lại để người kế tiếp trở thành hàng xóm của người thắng.
Cho ma trận kết quả quyết đấu. Nếu $A_{i,j}=1$ thì người $i$ luôn thắng người
$j$; nếu $A_{i,j}=0$ thì người $i$ thua người $j$. Người $k$ được xem là có
thể thắng nếu tồn tại một thứ tự rút thăm gồm $n-1$ trận để $k$ thắng trận
cuối.

Hãy đọc ma trận $A$ từ tệp \texttt{MUS.IN}, tính các số hiệu người có thể
thắng, rồi ghi chúng vào \texttt{MUS.OUT}.

\paragraph{Dữ liệu vào.}
Dòng đầu của \texttt{MUS.IN} chứa số nguyên $n$ với $3\le n\le 100$. Trong
mỗi dòng của $n$ dòng tiếp theo có một xâu gồm $n$ chữ số $0$ hoặc $1$. Chữ
số ở vị trí $j$ của dòng $i$ biểu diễn $A_{i,j}$. Với $i\ne j$, ta có
$A_{i,j}=1-A_{j,i}$, và giả sử $A_{i,i}=1$ với mọi $i$.

\paragraph{Dữ liệu ra.}
Dòng đầu của \texttt{MUS.OUT} ghi $m$, số người có thể thắng. Trong $m$ dòng
tiếp theo, ghi số hiệu của những người ấy theo thứ tự tăng dần, mỗi dòng một
số.

\paragraph{Ví dụ.}
\begin{verbatim}
Input
7
1111101
0101100
0111111
0001101
0000101
1101111
0100001

Output
3
1
3
6
\end{verbatim}
Một thứ tự quyết đấu giúp người số $6$ thắng là
$1$--$2$, $1$--$3$, $5$--$6$, $7$--$1$, $4$--$6$, $6$--$1$. Chỉ còn hai người
khác có thể thắng là $1$ và $3$.

\section{Bài Experiment}

Gần đây, các nhà khoa học đã dùng tàu thăm dò để lấy lần lượt $n_1$ và $n_2$
mảnh thiên thạch có cấu trúc tương tự từ sao $A$ và sao $B$. Cần thực hiện
$n=\min(n_1,n_2)$ thí nghiệm hóa học để so sánh chúng; mỗi thí nghiệm lấy một
mảnh từ sao $A$ và một mảnh từ sao $B$. Do biến đổi hóa học trong thí nghiệm,
một mảnh đã dùng thì không thể dùng lại. Vì đặc tính của loại thiên thạch này
bị ảnh hưởng lớn bởi khối lượng, các nhà khoa học muốn chọn một phương án sao
cho tổng trị tuyệt đối hiệu khối lượng của hai mảnh trong mỗi thí nghiệm là
nhỏ nhất.

Hãy tính tổng nhỏ nhất ấy.

\paragraph{Dữ liệu vào.}
Tệp vào là \texttt{Exp.dat}. Dòng đầu gồm hai số nguyên $n_1,n_2$. Tiếp theo
là $n_1$ dòng, mỗi dòng một số thực, biểu diễn khối lượng các mảnh từ sao
$A$. Sau đó là $n_2$ dòng, mỗi dòng một số thực, biểu diễn khối lượng các
mảnh từ sao $B$. Có $n_1,n_2\le 500$, và khối lượng mỗi mảnh không vượt quá
$90$.

\paragraph{Dữ liệu ra.}
Tệp ra là \texttt{Exp.out}, gồm một số thực duy nhất là tổng nhỏ nhất cần tìm.

\paragraph{Ví dụ.}
\begin{verbatim}
Input
4 2
44.9
50.0
77.2
86.4
59.8
58.9

Output
23.8
\end{verbatim}

\end{document}
